\documentclass[10pt,a4paper]{amsart}
\usepackage{amsmath,amssymb,amsthm,graphicx}
% Retained preamble for the Stacks Project chapter morphisms.tex (Morphisms of Schemes).
% The chapter file pinned-original-source.tex carries only \input{preamble} on its first line, so no
% preamble text lives inside the chapter itself.  The Stacks Project uses one single shared preamble for
% every chapter; the snapshot stages it verbatim in sources/stacks-chow/preamble.tex.  The block below is
% copied mechanically, line for line, from that file's lines 14-18 (the Xy-pic setup) and lines 154-217
% (the native theorem declarations together with the mathematical macros).  Nothing was added, removed or
% redefined.  The wrapper class and the amsmath/amssymb/amsthm/graphicx base packages are supplied by the
% assembler, so the class/\externaldocument part of the shared preamble is not repeated here.
% For commutative diagrams we use Xy-pic
\usepackage[all]{xy}

% We use 2cell for 2-commutative diagrams.
\xyoption{2cell}
% Theorem environments.
%
\theoremstyle{plain}
\newtheorem{theorem}[subsection]{Theorem}
\newtheorem{proposition}[subsection]{Proposition}
\newtheorem{lemma}[subsection]{Lemma}

\theoremstyle{definition}
\newtheorem{definition}[subsection]{Definition}
\newtheorem{example}[subsection]{Example}
\newtheorem{exercise}[subsection]{Exercise}
\newtheorem{situation}[subsection]{Situation}

\theoremstyle{remark}
\newtheorem{remark}[subsection]{Remark}
\newtheorem{remarks}[subsection]{Remarks}

\numberwithin{equation}{subsection}

% Macros
%
\def\lim{\mathop{\mathrm{lim}}\nolimits}
\def\colim{\mathop{\mathrm{colim}}\nolimits}
\def\Spec{\mathop{\mathrm{Spec}}}
\def\Hom{\mathop{\mathrm{Hom}}\nolimits}
\def\Ext{\mathop{\mathrm{Ext}}\nolimits}
\def\SheafHom{\mathop{\mathcal{H}\!\mathit{om}}\nolimits}
\def\SheafExt{\mathop{\mathcal{E}\!\mathit{xt}}\nolimits}
\def\Sch{\mathit{Sch}}
\def\Mor{\mathop{\mathrm{Mor}}\nolimits}
\def\Ob{\mathop{\mathrm{Ob}}\nolimits}
\def\Sh{\mathop{\mathit{Sh}}\nolimits}
\def\NL{\mathop{N\!L}\nolimits}
\def\CH{\mathop{\mathrm{CH}}\nolimits}
\def\proetale{{pro\text{-}\acute{e}tale}}
\def\etale{{\acute{e}tale}}
\def\QCoh{\mathit{QCoh}}
\def\Ker{\mathop{\mathrm{Ker}}}
\def\Im{\mathop{\mathrm{Im}}}
\def\Coker{\mathop{\mathrm{Coker}}}
\def\Coim{\mathop{\mathrm{Coim}}}

% Boxtimes
%
\DeclareMathSymbol{\boxtimes}{\mathbin}{AMSa}{"02}

%
% Macros for moduli stacks/spaces
%
\def\QCohstack{\mathcal{QC}\!\mathit{oh}}
\def\Cohstack{\mathcal{C}\!\mathit{oh}}
\def\Spacesstack{\mathcal{S}\!\mathit{paces}}
\def\Quotfunctor{\mathrm{Quot}}
\def\Hilbfunctor{\mathrm{Hilb}}
\def\Curvesstack{\mathcal{C}\!\mathit{urves}}
\def\Polarizedstack{\mathcal{P}\!\mathit{olarized}}
\def\Complexesstack{\mathcal{C}\!\mathit{omplexes}}
% \Pic is the operator that assigns to X its picard group, usage \Pic(X)
% \Picardstack_{X/B} denotes the Picard stack of X over B
% \Picardfunctor_{X/B} denotes the Picard functor of X over B
\def\Pic{\mathop{\mathrm{Pic}}\nolimits}
\def\Picardstack{\mathcal{P}\!\mathit{ic}}
\def\Picardfunctor{\mathrm{Pic}}
\def\Deformationcategory{\mathcal{D}\!\mathit{ef}}

% --- appended 2026-10-12 (Chief): the project own definitions of the reference/slogan/history
% display environments, copied mechanically from the shared Stacks preamble lines 8-12
% (\usepackage{verbatim} plus the three \newenvironment lines, which the chapter preamble keeps
% verbatim from the project preamble). Nothing added, removed or redefined: these environments
% carry no mathematical content and the project itself typesets them as comments. Without these
% three lines a counted statement whose own text uses them fails with Environment slogan undefined.
\usepackage{verbatim}
\newenvironment{reference}{\comment}{\endcomment}
\newenvironment{slogan}{\comment}{\endcomment}
\newenvironment{history}{\comment}{\endcomment}
\usepackage{cleveref}
\crefname{originallemma}{original lemma}{original lemmas}
\Crefname{originallemma}{Original lemma}{Original lemmas}
\makeatletter
\newlabel{properties-lemma-locally-constructible}{{054C}{source}{}{sourcepointer-1}{}}
\newlabel{properties-lemma-locally-constructible@cref}{{[originallemma][0][]054C}{[1][0][]source}{}{}{}}
\makeatother
\crefname{originallemma}{original lemma}{original lemmas}
\Crefname{originallemma}{Original lemma}{Original lemmas}
\makeatletter
\newlabel{schemes-lemma-quasi-compact-preserved-base-change}{{01K5}{source}{}{sourcepointer-2}{}}
\newlabel{schemes-lemma-quasi-compact-preserved-base-change@cref}{{[originallemma][0][]01K5}{[1][0][]source}{}{}{}}
\makeatother
\crefname{originallemma}{original lemma}{original lemmas}
\Crefname{originallemma}{Original lemma}{Original lemmas}
\makeatletter
\newlabel{topology-lemma-inverse-images-constructibles}{{005I}{source}{}{sourcepointer-3}{}}
\newlabel{topology-lemma-inverse-images-constructibles@cref}{{[originallemma][0][]005I}{[1][0][]source}{}{}{}}
\makeatother
\title{The inverse image of a constructible subset along any morphism is constructible}
\author{The Stacks Project Authors}
\date{stacks-project:morphisms.tex:local-snapshot-20261012; GNU FDL 1.2 or later}
\begin{document}
\maketitle
\noindent\textit{Editorial note.} All delivered mathematical text is the author's own, copied from the pinned Stacks Project chapter file pinned-original-source.tex; no statement or proof step was rewritten, reordered or removed.\par\medskip
\noindent\textit{Editorial note.} Following the multi-claim convention for this source, source\_url is the source-level URL https://stacks.math.columbia.edu/ and the source label lemma-inverse-image-constructible is carried in claim\_key and here rather than in a per-claim tag URL.\par\medskip
\noindent\textit{Editorial note.} The preamble is the chapter's shared Stacks Project preamble, retained as sources/stacks-morphisms/retained-preamble.tex: a line-for-line excerpt of the shared Stacks preamble staged with this snapshot. No macro was invented or redefined, and this card adds no preamble text of its own.\par\medskip
\noindent\textit{Editorial note.} Statement span 4339-4345 and proof span 4347-4366 come from the halfA pool card operations/pool-cards/stacks-morphisms-halfA.json and were re-read in the pinned chapter file before this card was written.\par\medskip
\noindent\textit{Editorial note.} The proof ends with the author's phrase 'we win'; it is preserved verbatim.\par\medskip
\noindent\textit{Editorial note.} 3 cross-references occur inside the counted spans; they are listed in external\_references as navigation pointers only.\par\medskip
\noindent\textit{Allowed standard prerequisites (editorial).}\begin{itemize}
\item The reduction of constructibility to retrocompact opens, cited from the topology chapter: it suffices to show that inverse images of retrocompact opens are retrocompact. Cited at proof lines 4564-4567. Reduces the constructible case to a statement about retrocompact opens.
\item An open immersion with retrocompact source is quasi-compact, and quasi-compactness survives base change, cited from the schemes chapter. Cited at proof lines 4568-4571. Shows that the base change of the retrocompact open is quasi-compact, which is exactly retrocompactness in the source.
\item The characterization of locally constructible subsets by their traces on affine opens, cited from the properties chapter. Cited at proof lines 4578-4580. Reduces the locally constructible case to the constructible case on an affine chart.
\end{itemize}\medskip
\ifdefined\hypertarget\hypertarget{sourcepointer-1}{}\fi
\noindent\textit{Source reference.} Original lemma 054C: focused-seeds/AGP/sources/stacks-snapshot-20261011/tags.txt, source lines 6635--6635. Same-chapter reference used in the retained author text at source line(s) 4362, 4415, 4433. It is supplied only as a navigation pointer so that the retained text compiles; the counted item delivers the author's own statement and the author's complete proof, and no part of the referenced item is copied in or substituted for the counted argument.\par\medskip
\ifdefined\hypertarget\hypertarget{sourcepointer-2}{}\fi
\noindent\textit{Source reference.} Original lemma 01K5: focused-seeds/AGP/sources/stacks-snapshot-20261011/tags.txt, source lines 2038--2038. Same-chapter reference used in the retained author text at source line(s) 1319, 2335, 2748, 3954, 4177, 4356. It is supplied only as a navigation pointer so that the retained text compiles; the counted item delivers the author's own statement and the author's complete proof, and no part of the referenced item is copied in or substituted for the counted argument.\par\medskip
\ifdefined\hypertarget\hypertarget{sourcepointer-3}{}\fi
\noindent\textit{Source reference.} Original lemma 005I: focused-seeds/AGP/sources/stacks-snapshot-20261011/tags.txt, source lines 210--210. Same-chapter reference used in the retained author text at source line(s) 4351. It is supplied only as a navigation pointer so that the retained text compiles; the counted item delivers the author's own statement and the author's complete proof, and no part of the referenced item is copied in or substituted for the counted argument.\par\medskip
\begin{lemma}
\label{lemma-inverse-image-constructible}
Let $f : X \to Y$ be a morphism of schemes.
Let $E \subset Y$ be a subset.
If $E$ is (locally) constructible in $Y$, then $f^{-1}(E)$ is (locally)
constructible in $X$.
\end{lemma}

\begin{proof}
To show that the inverse image of every constructible subset is constructible
it suffices to show that the inverse image of every retrocompact open $V$
of $Y$ is retrocompact in $X$, see
Topology, Lemma \ref{topology-lemma-inverse-images-constructibles}.
The significance of $V$ being retrocompact
in $Y$ is just that the open immersion $V \to Y$ is quasi-compact.
Hence the base change $f^{-1}(V) = X \times_Y V \to X$ is quasi-compact
too, see
Schemes, Lemma \ref{schemes-lemma-quasi-compact-preserved-base-change}.
Hence we see $f^{-1}(V)$ is retrocompact in $X$.
Suppose $E$ is locally constructible in $Y$.
Choose $x \in X$. Choose an affine neighbourhood $V$ of $f(x)$ and
an affine neighbourhood $U \subset X$ of $x$ such that $f(U) \subset V$.
Thus we think of $f|_U : U \to V$ as a morphism into $V$. By
Properties, Lemma \ref{properties-lemma-locally-constructible}
we see that $E \cap V$ is constructible in $V$. By the constructible case
we see that $(f|_U)^{-1}(E \cap V)$ is constructible in $U$.
Since $(f|_U)^{-1}(E \cap V) = f^{-1}(E) \cap U$ we win.
\end{proof}

\end{document}
